\chapter{Limits: What Is and Is Not Claimed}\label{ch:limits}

\begin{workplan}{\Recast}
\wpgoal{Keep this chapter's function --- the honest boundary --- and update
its content to V2's boundary, which is in some ways wider and in most ways
narrower than V0's.}

\wpsource{This chapter's existing text;
\texttt{MATH\_TOOLBOX.md}~Part~VIII (the consolidated dead-end registry),
Part~VI (four failed cover-detection attempts), Part~IX (the methodology
findings); \texttt{TATIANA\_V1\_IDENTITY\_MAP.md}~\S2, \S11, \S12.}

\wpsteps{
  \item \textbf{Replace the scope statement.} The V0 text scopes to three
    papers and a research-tool track. That track is retired. State V2's scope:
    one architecture, evaluated by whether its invariant survives its own
    dynamics.
  \item \textbf{Import the dead-end registry wholesale} into
    Appendix~\ref{app:refuted} and reference it here. Everything refuted,
    rejected or dissolved, with the reason, so none of it is re-proposed
    without knowing why it died. This is one of the project's genuine assets.
  \item \textbf{Keep the cover-vs-partition saga at full length.} Four
    instruments, all inconclusive, each diagnosed: an aggregation unit that
    fixed the answer before any data was seen; a null model with free row
    margins that manufactured a false positive on identical data ($p = 0.45$
    versus $p = 0.02$ purely by changing the null); an instrument blind
    without dimensionality reduction; a label/score relation mismatch. State
    the acceptance criterion earned from it: any instrument must pass a
    synthetic positive, a synthetic negative \emph{with planted structure
    recovered in the right direction}, and a structure-free geometric control,
    at one shared parameter setting, \emph{before} touching real data. None of
    the four met it.
  \item \textbf{Keep ``four failures to reject the partition are not evidence
    that the partition is true.''} Verbatim.
  \item \textbf{State the live external threat and the defence.} Published
    work reports that identity maps can replace learned sheaf Laplacians on a
    benchmark family. It concerns \emph{graphs}. The $2$-cell argument of
    Chapter~\ref{ch:complex} is the defence and must appear before a reader
    raises the objection --- and it must be stated as a theorem, since as of
    now it is asserted.
  \item \textbf{Reproduce the kill-tests as a live checklist}, with current
    status: routing (unrun), localisation on real data (unrun), invariant
    drift (\Measured, survived, with pruning untested), sheaf-not-decorative
    (unrun). Anything that fires goes in the logbook and in
    Appendix~\ref{app:refuted}.
  \item \textbf{Keep the prior-art honesty.} Cover detection has mature
    machinery in at least four fields under different names, and roughly a
    week was spent rediscovering it. What remains genuinely original is the
    sheaf over the structure, the growth law, the accumulating-not-recomputing
    memory, and the two-complex architecture. Infrastructure should be
    borrowed, not invented.
  \item \textbf{Carry the citation warning.} Nearly every citation in the
    source corpus is flagged as unsearched. Verify before any bibliography.
}

\wpdone{A reader finishing this chapter can state exactly what has been
established, what has been refuted, and what has merely not yet failed.}
\end{workplan}

\begin{intu}\Intu
The deepest thing the topology buys is not that it makes the system
smart. It is that it makes the claim of augmentation \emph{falsifiable}:
memory, consistency and growth become operations on one object, with
cheap provable guarantees and a live measure of whether augmentation is
actually happening. A framework that cannot fail is not doing any work,
and most of this chapter is about where this one can.
\end{intu}

\section{The bet and the refusal}

\begin{hon}\Hon
\textbf{The bet (defensible).} A persistent, self-checking,
self-restructuring \emph{organiser} wrapped around a fixed language
model can, over time, produce better-grounded and more internally
consistent research artifacts than the raw model, through measured
cross-organ error correction and accumulation. The advantage is
\emph{architectural}, expressed as operations on the single object
$\M=(\CC,\Fs,s)$.

\textbf{The refusal.} No simplex proves a theorem and no Laplacian
integrates a differential form. The per-inference reasoning quality is
bounded by the underlying model. The mathematics here does not raise
that ceiling; it organises, measures and error-corrects \emph{around}
it.
\end{hon}

The bet is about the \emph{slope}, not the intercept. That makes it an
empirical claim, and Open Problem~\ref{op:slope} records that it has not
been tested.

\section{The scope is three papers, not one}

\begin{hon}\Hon
A single document containing the $\rhoscore$ splitting, the operator
algebra, the coned concept, the cohomological obstruction and a spectral
prediction on $\mathbb{CP}^3$ will produce something no referee can
accept, because those claims have incompatible standards of evidence.
The audit's recommended split is:

\textbf{Paper A --- a cellular-sheaf coherence measure for multi-agent
orchestration.} Content: $\M=(\CC,\Fs,s)$; $\discord$ and $\rhoscore$;
the $\dissent/\rhotil$ splitting with the window; weighted
Anderson--Morley \emph{with} the $\norm{\rest{}}\le1$ hypothesis stated;
worst-edge localisation; the three-valued verification gate. Every claim
is a theorem plus a measurement on a running system. This paper is
nearly writable today.

\textbf{Paper B --- sheaf-cohomological obstruction as a diagnostic.}
Content: Proposition~\ref{prop:vacuous} --- the derived cochain is exact,
a negative result worth publishing on its own; the architecture needed
to produce non-derived $1$-cochains; the Hodge trichotomy; and the
connection to contextuality in the Abramsky--Brandenburger sense, which
is the thread that would make this a quantum-information contribution
rather than a software paper. This paper \emph{requires new engineering}
and cannot be written before that lands.

\textbf{Paper C --- the spectral work.} This is separate mathematics. It
has nothing to do with \textsc{tatiana} and should not be contaminated
by it; the architecture's role is to be the tool that helped, mentioned
in a methods note, not a co-author of the theorem.
\end{hon}

\section{What must not be claimed}

Carrying forward the honesty discipline that is the best thing about the
source corpus:

\begin{itemize}
  \item \textbf{Not} ``we prove the whole exceeds the sum''. A
        structural novelty statement says a composite is \emph{new}, not
        that it is true or useful. Verification remains the sole
        arbiter and runs first.
  \item \textbf{Not} ``capacity measures learning''. It is gameable by
        adding independent useless operators. Report it; never optimise
        it.
  \item \textbf{Not} ``the system uses $D$-module theory''. It uses
        Gaussians, which \emph{are} the simplest holonomic modules --- an
        exact identification that currently buys three invariants, all
        of them constant. The defensible claim is that the concept type
        admits a conservative extension to holonomic modules, feasible
        only because of the existing low-rank design.
  \item \textbf{Not} ``Riemann--Hilbert applies''. One leg is rigorous
        (cellular sheaves are constructible sheaves via exit paths); the
        other requires a commitment to an algebraic stratification that
        has not been made.
  \item \textbf{Not} ``the threshold is derived''. It is
        \emph{normalised} (Proposition~\ref{prop:window}).
  \item \textbf{Not} ``$b_1$ measures conceptual gaps'' without
        Remark~\ref{rem:b1earned}. The justification is the whole
        content of the claim.
  \item \textbf{Not} ``the architecture is an operad''. See
        Remark~\ref{rem:tracemonoid} and the three defects following it.
\end{itemize}

\section{The findings, restated as boundaries}

The audit's fourteen findings are not a list of bugs; each marks a place
where a claim outruns its evidence. Restated as standing limits:

\begin{enumerate}
  \item \textbf{The audit itself was against a stale snapshot.} Several
        findings were repaired before the audit was read. Any use of
        this list must be re-checked against the current branch --- and
        that instruction applies to this book.
  \item \textbf{Theorem 7.5 is stated more generally than proved.} The
        gap between hypothesis and proof has not been located
        precisely, and until it is, the theorem should be cited in its
        proved form only.
  \item \textbf{The $\Hone$ repair is half circular.} The circularity
        must be identified before the repair is adopted
        (Open Problem~\ref{op:H1}).
  \item \textbf{Proposition 12.4 is false in the deployed case.} A
        counterexample exists and must be reproduced in
        Appendix~\ref{app:refuted} before that entry is complete.
  \item \textbf{The $K_0$ classes are trivial for every object the
        system produces.} An invariant that separates nothing is not a
        diagnostic (Open Problem~\ref{op:jordanholder}).
  \item \textbf{The independence relation is not symmetric}, its
        supports are placeholders, and its empty-set convention is
        inverted (Chapter~\ref{ch:operad}).
  \item \textbf{The restriction maps are the identity}, so the sheaf is
        a graph Laplacian in disguise
        (Open Problem~\ref{op:restrictions}).
\end{enumerate}

\begin{hon}\Hon
The honest summary of the whole corpus: the \emph{measurement} layer is
real, proved and running --- $\discord$, $\rhoscore$, the splitting, the
window, the three-valued gate, the coned concept and the growth law all
survive scrutiny. The \emph{expressive} layer is not yet built: while
restrictions are the identity and $1$-cochains are derived, the sheaf
and its cohomology are carrying no information the graph did not already
carry. That is not a fatal criticism; it is a precise statement of what
the next piece of work has to be, and the fact that it can be stated
this precisely is what the formalism bought.
\end{hon}
